\documentclass[12pt,a4paper]{article}
\usepackage[utf8]{inputenc}
\usepackage[french]{babel}
\usepackage[T1]{fontenc}
\usepackage{amsmath,amssymb}
\usepackage{graphicx}
\usepackage{xcolor}
\usepackage{hyperref}
\usepackage{geometry}
\usepackage{tcolorbox}
\usepackage{tikz}
\usepackage{array}
\geometry{margin=2.5cm}
% Définition des couleurs
\definecolor{titrebleu}{RGB}{0,102,204}
\definecolor{defvert}{RGB}{0,153,76}
\definecolor{exempleorange}{RGB}{255,102,0}
\definecolor{algoviolet}{RGB}{153,0,153}
\definecolor{importantrouge}{RGB}{204,0,0}
\definecolor{fondjaune}{RGB}{255,250,205}
\definecolor{fondvert}{RGB}{230,255,230}
\definecolor{fondbleu}{RGB}{230,240,255}
% Boîtes colorées
\newtcolorbox{definition}{
    colback=fondvert,
    colframe=defvert,
    fonttitle=\bfseries,
    title=�� Définition,
    arc=3mm
}
\newtcolorbox{exemple}{
    colback=fondjaune,
    colframe=exempleorange,
    fonttitle=\bfseries,
    title=�� Exemple,
    arc=3mm
}
\newtcolorbox{methode}{
    colback=fondbleu,
    colframe=algoviolet,
    fonttitle=\bfseries,
    title=⚙️ Méthode,
    arc=3mm
}
\newtcolorbox{remarque}{
    colback=white,
    colframe=importantrouge,
    fonttitle=\bfseries,
    title=�� Remarque,
    arc=3mm
}
% Style des sections
\usepackage{titlesec}
\titleformat{\section}
{\color{titrebleu}\normalfont\LARGE\bfseries}
{\color{titrebleu}\thesection}{1em}{}
\titleformat{\subsection}
{\color{defvert}\normalfont\Large\bfseries}
{\color{defvert}\thesubsection}{1em}{}
\titleformat{\subsubsection}
{\color{algoviolet}\normalfont\large\bfseries}
{\color{algoviolet}\thesubsubsection}{1em}{}
\title{\textbf{\Huge\color{titrebleu}Fiche de Cours}\\[0.5em]
\large\color{defvert}Architecture des Systèmes\\
Chapitre 1 : Fonctionnement d'un ordinateur}
\author{\large Alexandre Guitton}
\date{\large Version 2024-2025}
\begin{document}
\maketitle
\tableofcontents
\newpage

\section{Objectifs du Chapitre}
\begin{definition}
\begin{itemize}
    \item Fonctionnement d'un ordinateur
    \item Stockage : des bits aux valeurs
    \item Traitement :
    \begin{itemize}
        \item Par des circuits simples (sans notion de temps et sans stockage)
        \item Par des circuits complexes (avec notions de temps et de stockage)
        \item Entre circuits complexes
    \end{itemize}
\end{itemize}
\end{definition}

\section{Plan du Chapitre}
\begin{itemize}
    \item 1. Représentation des valeurs
    \item 2. Circuits logiques
    \item 3. Circuits séquentiels
    \item 4. Architecture d'un ordinateur
\end{itemize}

\newpage
\section{Représentation des valeurs}

\subsection{Bit et Stockage}
\begin{definition}
\textbf{Bit} : Unité de stockage de l'information, peut être 0 ou 1.
\end{definition}

\begin{remarque}
Stockage de valeurs complexes comme nombres, textes, images, sons, vidéos via :
\begin{itemize}
    \item Entiers positifs, entiers naturels
    \item Nombres réels
    \item Caractères, chaînes de caractères
\end{itemize}
\end{remarque}

\subsection{Entiers positifs}

\subsubsection{Bases}
\begin{definition}
\textbf{Base} : Nombre de chiffres différents autorisés.
\begin{itemize}
    \item Base 2 (binaire) : 0 et 1
    \item Base 10 (décimale) : 0 à 9
    \item Base 16 (hexadécimale) : 0 à 9, A à F
\end{itemize}
\end{definition}

\begin{remarque}
Notation : Base en indice, hexadécimal précédé de 0x. Par défaut, base 10.
\end{remarque}

\subsubsection{Conversion base b vers base 10}
\begin{methode}
Pour $(a_{n-1} \dots a_0)_b$ : $\sum_{i=0}^{n-1} a_i \cdot b^i$
\end{methode}

\begin{exemple}
$110101_2 = 1\cdot2^5 + 1\cdot2^4 + 0\cdot2^3 + 1\cdot2^2 + 0\cdot2^1 + 1\cdot2^0 = 53_{10}$
\end{exemple}

\subsubsection{Conversion base 10 vers base b}
\begin{methode}
\begin{enumerate}
    \item Écrire puissances de b
    \item Trouver c (0 à b-1) et k tels que $c \cdot b^k \leq n < (c+1) \cdot b^k$
    \item Soustraire $c \cdot b^k$, répéter jusqu'à 0
\end{enumerate}
\end{methode}

\begin{exemple}
39 en base 2 : 100111$_2$
\end{exemple}

\subsubsection{Conversions entre bases}
\begin{methode}
Passer via base 10. Pour base 16 ↔ base 2 : 4 bits par chiffre hex.
\end{methode}

\subsubsection{Addition}
\begin{methode}
Comme en décimal, avec retenues. Attention aux débordements.
\end{methode}

\begin{exemple}
$100111_2 + 101110_2 = 1010101_2$ (39 + 46 = 85)
\end{exemple}

\newpage
\subsection{Entiers naturels}

\begin{remarque}
Problème : Représenter nombres négatifs.
\end{remarque}

\subsubsection{Bit de signe}
\begin{exemple}
39 sur 8 bits : 00100111 ; -39 : 10100111
\end{exemple}

\begin{remarque}
Problèmes : Additions/soustractions différentes, 0 et -0 distincts.
\end{remarque}

\subsubsection{Complément logique (à 1)}
\begin{methode}
Inverser bits pour négatifs.
\end{methode}

\begin{exemple}
-39 sur 8 bits : 11011000
\end{exemple}

\begin{methode}
Addition : Addition base 2, réajouter retenue débordante.
\end{methode}

\begin{remarque}
Problèmes : Addition complexe, 0 et -0 distincts.
\end{remarque}

\subsubsection{Complément arithmétique (à 2)}
\begin{methode}
Ajouter 1 au complément logique pour négatifs.
\end{methode}

\begin{exemple}
-39 sur 8 bits : 11011001
\end{exemple}

\begin{methode}
Addition : Addition base 2, ignorer retenue débordante.
\end{methode}

\newpage
\subsection{Nombres réels}

\begin{definition}
Format IEEE 754 (simple/double précision) :
\begin{itemize}
    \item Signe : 1 bit
    \item Mantisse normalisée : $m \cdot 2^e$, $0.5 \leq m < 1$
    \item Bit caché : Premier bit (1) non stocké
    \item Exposant biaisé
\end{itemize}
\end{definition}

\begin{remarque}
Simple : 1 signe, 23 mantisse, 8 exposant (biais -127). Double : 1, 52, 11 (-1023).
\end{remarque}

\begin{exemple}
11000000 00110000 00000000 00000000 = -2.75
\end{exemple}

\begin{exemple}
4.25 = 01000000 10001000 00000000 00000000$_2$
\end{exemple}

\subsection{Caractères}

\begin{definition}
Encodages : ASCII, UTF-8, etc.
\end{definition}

\begin{remarque}
ASCII : 128 caractères sur 1 octet. Ex : 'A' = 65.
\end{remarque}

\begin{remarque}
Propriétés : Suites pour minuscules/majuscules/chiffres. Problèmes : Pas d'accents, caractères non affichables.
\end{remarque}

\subsection{Chaînes de caractères}
\begin{methode}
\begin{itemize}
    \item Longueur + liste (Pascal)
    \item Caractère terminal (C)
\end{itemize}
\end{methode}

\newpage
\section{Circuits logiques}

\begin{definition}
Circuit électronique calculant sortie binaire en fonction d'entrées binaires.
\end{definition}

\begin{definition}
Fonctions : NOT, OR, AND, XOR, NOR, NAND.
\end{definition}

\subsection{Algèbre de Boole}
\begin{definition}
Opérations : NOT ($\overline{A}$), OR (+), AND ($\cdot$).
\end{definition}

\begin{remarque}
Lois de De Morgan : $\overline{A+B} = \overline{A} \cdot \overline{B}$, etc.
\end{remarque}

\subsection{Table de vérité}
\begin{methode}
Liste exhaustive entrées/sorties.
\end{methode}

\begin{exemple}
L allumée si A ou B ouverts, sauf si S ouvert : L = $\overline{S} \cdot (A + B)$
\end{exemple}

\subsection{Synthèse}
\begin{methode}
Table de vérité → expression → circuit.
\end{methode}

\subsection{Analyse}
\begin{methode}
Circuit → expression.
\end{methode}

\subsection{Simplification}
\begin{methode}
Observation, propriétés Boole, Karnaugh.
\end{methode}

\newpage
\section{Circuits séquentiels}

\begin{definition}
Intègrent temps (horloge) et mémoire (bascule).
\end{definition}

\subsection{Bascule RS}
\begin{definition}
Entrées R (reset à 0), S (set à 1).
\end{definition}

\begin{methode}
Table : Si R=1 S=0 → 0 ; R=0 S=1 → 1 ; R=0 S=0 → maintient ; R=1 S=1 → indéterminé.
\end{methode}

\begin{definition}
$Q_{new} = \overline{R} \cdot (S + Q_{old})$
\end{definition}

\newpage
\section{Architecture d'un ordinateur}

\begin{definition}
Interconnexion circuits séquentiels + horloge + bus. Modèle Von Neumann.
\end{definition}

\subsection{Types de circuits}
\begin{itemize}
    \item Mémoire
    \item Unités de calculs
    \item Unités de traitement
    \item Périphériques externes
\end{itemize}

\subsection{Mots-clés}
\begin{definition}
\begin{itemize}
    \item RAM : Mémoire principale lecture/écriture
    \item ROM : Lecture seule
    \item BIOS : Primitives E/S en ROM
    \item CMOS : Sauvegarde BIOS éteint
    \item Circuit intégré : Fonctions sur puce
    \item Processeur : Traitements programmables
    \item Microprocesseur : Processeur sur un circuit intégré
\end{itemize}
\end{definition}

\section{Résumé des connaissances}
\begin{itemize}
    \item Codage : Entiers positifs/naturels, réels, caractères
    \item Circuits logiques : Synthèse, analyse
    \item Circuits séquentiels : Principe
    \item Ordinateur : Organisation
\end{itemize}

\end{document}